\documentclass[11pt]{article}
\usepackage[margin=1in]{geometry}
\usepackage{amsmath,amssymb,amsthm,booktabs,array,enumitem,tabularx,hyperref}
\newtheorem{definition}{Definition}
\newtheorem{proposition}{Proposition}
\newcommand{\IR}{\mathsf{IR}}
\newcommand{\FE}{\mathsf{FE}}
\newcommand{\GE}{\mathsf{GE}}
\newcommand{\EN}{\mathsf{EN}}
\title{Four Project-Success Categories and Their Evidence Obligations\\\large P14-T01 task-local draft/control formalization}
\author{The AEther-Flow Research Project}
\date{22 July 2026}

\begin{document}
\maketitle

\begin{abstract}
This task-local packet separates interpretive redescription, formal or
categorical equivalence, genuine emergence, and empirical novelty.  Each
category has its own evidence certificate.  A finite countermodel theorem
shows that none implies another by default.  The classification is
proposal-neutral: it does not edit ontology, promote a physical
interpretation, change the Distance-to-GR ledger, or establish a completed
derivation.
\end{abstract}

\section{Scope and status algebra}
Let $x$ be a precisely scoped project claim and let $E(x)$ be the tracked
evidence available for that claim.  Define four predicates
\[
  \IR(x),\quad \FE(x),\quad \GE(x),\quad \EN(x),
\]
and the tri-valued status vector
\[
  \Sigma(x)=(s_{\IR},s_{\FE},s_{\GE},s_{\EN})
  \in\{\texttt{met},\texttt{not\_met},\texttt{indeterminate}\}^{4}.
\]
Here \texttt{not\_met} means that the declared evidence obligation is not
satisfied; it is not an impossibility theorem.  \texttt{indeterminate} means
that the scoped record does not currently decide the obligation.  No status
component is copied into another component.

\section{Four independent evidence certificates}

\begin{definition}[Interpretive redescription, \texorpdfstring{$\IR$}{IR}]
A claim satisfies \emph{interpretive redescription} on an operational domain
$D$ only when it supplies (i) an explicit dictionary between the benchmark
formalism and the proposed interpretive vocabulary, (ii) a declared domain
and approximation regime, and (iii) evidence that the mapped description
preserves the benchmark's operational predictions on $D$.  The certificate
may explain the same successful formalism differently; it does not thereby
establish ontology adoption, categorical equivalence, emergence, or empirical
novelty.
\end{definition}

Necessary evidence is a typed dictionary, a scope statement, and an
operational-preservation witness.  Merely renaming variables, drawing an
analogy, matching a few equations, or asserting that the interpretation is
intuitive is insufficient.

\begin{definition}[Formal or categorical equivalence, \texorpdfstring{$\FE$}{FE}]
A claim satisfies \emph{formal or categorical equivalence} only relative to
declared formal systems or categories $\mathcal C$ and $\mathcal D$, a
preservation profile $\Pi$, and hypotheses $H$.  Its certificate contains a
proved isomorphism or an equivalence of categories: functors
$F\colon\mathcal C\to\mathcal D$ and $G\colon\mathcal D\to\mathcal C$
with the required natural isomorphisms, together with a proof that the
structures named in $\Pi$ are preserved (and reflected when claimed).
\end{definition}

Necessary evidence is an exact theorem identity, hypotheses, maps or functors,
proof obligations, and the preservation/reflection profile.  Similar
equations, equal parameter counts, a suggestive quotient, numerical agreement,
or shared terminology is insufficient.  A quotient-EqSrc construction enters
this category only when its declared equivalence theorem is actually proved.

\begin{definition}[Genuine emergence, \texorpdfstring{$\GE$}{GE}]
A claim satisfies \emph{genuine emergence} only when effective target
structure or semantics are derived from source primitives by a non-circular
reconstruction, coarse-graining, or limiting map.  The certificate names the
source signature, target object, bridge map, regime or scale, robustness
conditions, and a target-independence audit showing that the target structure
was not assumed in the source data.
\end{definition}

Necessary evidence is a source-to-effective construction plus regime,
robustness, and non-circularity evidence.  An interpretive dictionary, a
target embedding, a hand-chosen target atlas, formal equivalence alone, or
ontology adoption by stipulation is insufficient.

\begin{definition}[Empirical novelty, \texorpdfstring{$\EN$}{EN}]
A claim satisfies \emph{empirical novelty} only when it identifies a distinct,
falsifiable operational prediction or constraint not already fixed by the
declared benchmark.  Its certificate names an observable, test domain,
benchmark prediction set, project prediction set, a nonempty distinguishing
region or constraint, and uncertainty or identifiability conditions.
\end{definition}

Necessary evidence is an operational discriminator and a testable distinction.
Exact benchmark recovery, parameter relabeling, a new visualisation, a formal
equivalence, or a metaphysical preference is insufficient.

\section{Default non-implication theorem}

\begin{proposition}[SC4-NONIMPLICATION-THEOREM-V1]
Under the evidence-certificate definitions above, and absent an additional
task-specific implication theorem, no one of $\IR,\FE,\GE,\EN$ implies any
other.  Thus for every ordered pair $A\ne B$ there is a finite admissible
evidence world in which $A$ is met and $B$ is not met.
\end{proposition}

\begin{proof}
Consider four one-hot evidence worlds.  In $w_{\IR}$ an exact benchmark
reformulation has a typed interpretive dictionary and operational preservation,
but no categorical theorem, source reconstruction, or distinct prediction.
In $w_{\FE}$ two abstract formal categories have a proved structure-preserving
equivalence, but no benchmark interpretation, source-to-effective derivation,
or empirical content.  In $w_{\GE}$ a many-to-one source model robustly
reconstructs an effective target, but is neither an interpretive rewriting nor
an equivalence and produces no benchmark-distinguishing prediction.  In
$w_{\EN}$ a phenomenological rule gives a distinct falsifiable prediction but
supplies no interpretive dictionary, formal equivalence, or source derivation.
The truth vectors are respectively
\[
 (1,0,0,0),\quad(0,1,0,0),\quad(0,0,1,0),\quad(0,0,0,1).
\]
For any ordered pair $A\ne B$, the world whose sole met component is $A$
is a countermodel to $A\Rightarrow B$.  All twelve directed implications
therefore fail by default.  A later scoped theorem may prove an implication
under additional hypotheses, but it must name and discharge those hypotheses.
\end{proof}

The proposition is a task-local claim-logic result, not a theorem about nature
and not proof that any future source extension is impossible.

\section{Current project mapping}

\begin{center}
\small
\begin{tabularx}{\textwidth}{>{\raggedright\arraybackslash}p{0.20\linewidth}cccc>{\raggedright\arraybackslash}X}
\toprule
Claim instance & $\IR$ & $\FE$ & $\GE$ & $\EN$ & Reason \\
\midrule
Current exact-GR benchmark line & met & indeterminate & not met & not met & The work redescribes an operationally exact-GR benchmark; no source-to-GR formal-equivalence theorem is presently established, and the line supplies neither source-derived emergence nor a distinct prediction. \\
General quotient-EqSrc equivalence & indeterminate & indeterminate & not met & not met & A formal-equivalence status is conditional on a proved quotient/functor theorem with its preservation profile; scoped candidates do not discharge the general burden. \\
\bottomrule
\end{tabularx}
\end{center}

The first row is limited to benchmark-level interpretive redescription.  Its
formal-equivalence coordinate is indeterminate, not refuted: no scoped
source-to-GR equivalence certificate presently decides it.  The row
does not establish a physical gauge interpretation, a source law, $M_{\rm
src}$, $g_{\rm eff}$, matter coupling, Einstein-equation derivation, or
benchmark promotion.  The second row is deliberately conditional: a future
proof changes only the scoped $\FE$ component it actually establishes.

\section{Authority boundary}
These definitions and the non-implication proposition are `draft/control'.
They add no ontology or physical assumption.  They do not modify canonical
ontology, adopt or reject a source extension, infer theorem truth from a
validator, change scientific ledgers, authorize publication, or claim a
completed first-principles derivation of general relativity.

\end{document}
